\chapter{Green Revolution and Its Impacts}

\section{The Green Revolution and the Agrarian Revolution}

\begin{KeyConcept}
\begin{itemize}
\item \textbf{Revolution} = a change in the \textbf{social structure} (not mere reform); the \textbf{industrial revolution} moved work from the \textbf{private sphere (household)} to the \textbf{public sphere} and created the proletariat and bourgeoisie. \textbf{Green Revolution} was a type of \textbf{agrarian revolution} because it, too, changed the social structure --- new classes, tenants becoming labourers, women domesticated, the rise of OBC politics in the North and dominant-caste politics in the South, and even new religious movements (the \textbf{Dera gurus} --- e.g. Ram Rahim, Ram Pal). Earlier agriculture was a \textbf{household economy --- 'oikonomia'} --- from which we moved to the modern \textbf{'economy'}.
\item \textbf{Green Revolution} = the increase in the \textbf{productivity of agriculture} in India, based on \textbf{High-Yielding Variety (HYV) seeds} (developed by \textbf{Norman Borlaug}), \textbf{irrigation, mechanisation and chemicals} (fertilisers, pesticides) --- primarily of wheat and rice. It began in the period of \textbf{Lal Bahadur Shastri}, with \textbf{M. S. Swaminathan} as a key figure.
\item It succeeded in \textbf{Punjab, Haryana and Western UP} (later called the \textbf{'Bermuda Triangle of Girls'} because female feticide/infanticide began there), plus parts of Tamil Nadu, Kerala, Karnataka and Andhra --- about \textbf{one-third of India's arable area}.
\item The colonial government had used technology from the 1930s, but only for export cash-crops (tea, coffee, rubber); its wheat/rice seed experiments failed.
\end{itemize}
\end{KeyConcept}

\begin{QuickRevision}
\begin{itemize}
\item Revolution
\item social structure
\item industrial revolution
\item private sphere (household)
\item public sphere
\item Green Revolution
\item agrarian revolution
\item Dera gurus
\item household economy --- 'oikonomia'
\item 'economy'
\item productivity of agriculture
\item High-Yielding Variety (HYV) seeds
\item Norman Borlaug
\item irrigation, mechanisation and chemicals
\end{itemize}
\end{QuickRevision}

\section{Land Reform and Technological Intervention}

\begin{KeyConcept}
\begin{itemize}
\item Green Revolution was the \textbf{second stage of the land-reform programme}. Stage 1 = \textbf{redistribution of land} (land-ceiling acts); stage 2 = \textbf{technological intervention}. Both remain incomplete.
\item Land reform gave tenancy rights (not ownership), and its failure produced the \textbf{Naxalite movement} (Kanu Sanyal, Charu Majumdar). West Bengal and Kerala had the best land-reform laws, but in WB the reform was \textbf{successful only on paper} --- the previous landlords beat the new owners until the Naxalite movement forced possession.
\item The redistribution-vs-technology debate: Nehru feared technology's benefits would be captured by the upper class, but the USA (large plots) vs Japan (small plots, \emph{higher} productivity) comparison showed \textbf{productivity does not depend on farm size}. Land was pursued first because it carries human dignity and subsistence (India's history of famines is why migrants never sell their village land).
\item \textbf{Technological intervention} = \textbf{mechanisation} (high capital) + \textbf{genetically engineered seeds + fertilisers} (popular in the West after the 1930s; India had to wait for seeds suited to its climate, so Green Revolution came in the 1960s).
\item Two institutions powered it: the \textbf{developmental bureaucracy} and the \textbf{cooperatives}.
\end{itemize}
\end{KeyConcept}

\begin{QuickRevision}
\begin{itemize}
\item Revolution = change in social structure (work private$\rightarrow$public; new classes); Green Revolution = agrarian revolution.
\item GR = HYV seeds (Borlaug) + irrigation + mechanisation + chemicals; Shastri/Swaminathan; Punjab/Haryana/WUP ('Bermuda Triangle of Girls') + South; $\sim$1/3 arable area.
\item Land reform: redistribution (Naxalite movement --- Sanyal/Majumdar) $\rightarrow$ technological intervention; USA vs Japan (productivity not farm-size dependent).
\item Two institutions: developmental bureaucracy + cooperatives.
\end{itemize}
\end{QuickRevision}

\section{The Developmental Bureaucracy: Joan Mencher}

\begin{KeyConcept}
\begin{itemize}
\item \textbf{Joan Mencher} (fieldwork in \textbf{Chingleput}, Tamil Nadu) studied the role of the \textbf{developmental bureaucracy} in the Green Revolution. Every input (seeds from the National Seed Corporation, fertilisers, machine loans) was regulated by the state.
\item She found the bureaucracy was \textbf{not neutral}: it defined \textbf{small and marginal farmers} (less than 5 acres of irrigated land) as 'uneconomic', and regarded only the \textbf{big farmers as 'progressive farmers'} (investing in a capital-intensive agriculture). But the \textbf{small and marginal farmers were equally innovative and enthusiastic} about the new technology --- they simply were \textbf{not given the means} (cooperative loans, seeds, encouragement).
\end{itemize}
\end{KeyConcept}

\begin{QuickRevision}
\begin{itemize}
\item Joan Mencher
\item Chingleput
\item developmental bureaucracy
\item not neutral
\item small and marginal farmers
\item big farmers as 'progressive farmers'
\item small and marginal farmers were equally innovative and enthusiastic
\item not given the means
\end{itemize}
\end{QuickRevision}

\section{Cooperatives and Seeds: Satyadeva}

\begin{KeyConcept}
\begin{itemize}
\item \textbf{Satyadeva} studied the \textbf{Haryana Seeds Development Corporation} (part of the \textbf{National Seed Project}, started with the \textbf{World Bank}). To get seeds you had to buy a \textbf{share} (\rupee{}100 per acre for 5 crops).
\item Among his sample of \textbf{285 shareholder seed growers}: one-third grew seeds on $\sim$5 acres, one-third on 5--10 acres, and one-third on 10--100 acres --- so the policy \textbf{mainly benefited big farmers} (who had the money to pay, and who \textbf{controlled the cooperatives}, getting loans at $\sim$6\% against the bank's $\sim$12\%).
\item Most of the 285 shareholders were \textbf{upper-caste/dominant-caste}, and none were SC. The \textbf{public-sector seed corporations were in effect created to benefit the private sector} --- an expansion of the private sector in agriculture that, historically, always benefits the large farmers.
\end{itemize}
\end{KeyConcept}

\begin{QuickRevision}
\begin{itemize}
\item Mencher (Chingleput): bureaucracy biased --- 'small/marginal' = uneconomic; big = 'progressive'; but small farmers equally innovative, just not given means.
\item Satyadeva (Haryana): seed share \rupee{}100/acre; 285 growers (1/3 each: 5 / 5--10 / 10--100 acres); benefits big farmers (co-op loans 6\% vs bank 12\%); upper/dominant caste; PSUs benefiting private sector.
\end{itemize}
\end{QuickRevision}

\section{Bhalla and Chadda: The Pro-Rich Bias Fallacy}

\begin{KeyConcept}
\begin{itemize}
\item \textbf{Bhalla and Chadda} (agricultural economists) tested the hypothesis that Green Revolution technology, being \textbf{capital-intensive}, suits big farms better. Their Punjab study (3 regions, \textbf{180 villages} using the 30th round of the NSS methodology, \textbf{1,663 cultivating households}) \textbf{exposed the fallacy}.
\item Findings: \textbf{material cost per acre is inversely related to farm size} (small/marginal farmers spent \textbf{48.63\% of material cost on drought cattle}; rent also fell as farm size rose); \textbf{diesel/repair costs rose with farm size}. So the technology is \textbf{as efficient, if not more, on small farms}.
\item The 140 marginal farms (8.4\% of surveyed farms, avg. 1.6 acres) contrasted with large farms (avg. 32.8 acres --- about 20 times larger); \textbf{farm business income per acre was \rupee{}754.50 for the smallest farms vs \rupee{}740.50 for the large} --- small farmers are at least as productive, if not more. Such inequalities are then \textbf{perpetuated by the developmental bureaucracy's preference for large farms}.
\end{itemize}
\end{KeyConcept}

\begin{QuickRevision}
\begin{itemize}
\item Bhalla \& Chadda (Punjab, 180 villages, 1663 households): material cost inversely related to farm size; small farms as productive (\rupee{}754.50 vs \rupee{}740.50 per acre); pro-rich bias = fallacy; inequality perpetuated by bureaucracy.
\end{itemize}
\end{QuickRevision}

\section{Wet vs Dry Areas: Bagchi and Athreya}

\begin{KeyConcept}
\begin{itemize}
\item \textbf{Amiya Bagchi}: because the bureaucracy insisted Green Revolution suited \textbf{irrigated (wet) areas}, \textbf{inequality increased more in wet areas} (they became more prosperous, but with glaring inequality).
\item \textbf{B. Athreya} (Trichy district, Tamil Nadu; 6 villages from 2 panchayats): in \textbf{wet areas} poor peasants' agricultural income/wages (\rupee{}346) were \textbf{4.2 times higher} than in dry areas; agricultural labourers' cash wages were \rupee{}1,030 (wet) vs \rupee{}530 (dry) --- and \textbf{total income was 1.5 times higher} in wet areas. So \textbf{the economic condition of wage labourers was always better in wet areas}.
\item A striking finding: the \textbf{income of agricultural labourers was outstripping the peasants} (labourers \rupee{}1,686 vs peasants \rupee{}1,333) --- because labourers were on cash wages, while peasants (on crop-share) bore the rising input cost. So a \textbf{class structure formed}: landowner $\rightarrow$ tenant $\rightarrow$ peasant $\rightarrow$ labourer, with \textbf{peasants becoming poorer}.
\item \textbf{Peasant vs tenant}: the tenant does not work in the field (peasants work under him); the peasant works/hires labourers. Pre-Green-Revolution movements were \textbf{peasant movements}; the new farmers' movements are movements of \textbf{big landlords/tenants}.
\item \textbf{Mahalwari} (North India): village elders divided the common land, keeping the best (arable) land for themselves --- for the first time creating a class of \textbf{marginal farmers}.
\end{itemize}
\end{KeyConcept}

\begin{QuickRevision}
\begin{itemize}
\item Bagchi: inequality rose more in wet (irrigated) areas.
\item Athreya (Trichy): wet vs dry --- 4.2$\times$ (peasants), 2$\times$ (labourer wages), 1.5$\times$ (total income); labourers (\rupee{}1,686) outstripped peasants (\rupee{}1,333); class structure landlord$\rightarrow$tenant$\rightarrow$peasant$\rightarrow$labourer.
\item Peasant (works) vs tenant (doesn't); pre-GR peasant movements $\rightarrow$ post-GR big-farmer movements; Mahalwari $\rightarrow$ marginal farmers.
\end{itemize}
\end{QuickRevision}

\section{Scarlett Epstein: A Longitudinal Study (Mandya)}

\begin{KeyConcept}
\begin{itemize}
\item \textbf{Scarlett Epstein} did a \textbf{longitudinal study} (comparing the same area in 1955 and 1970) in \textbf{Mandya district, Karnataka}. After the Green Revolution's irrigation reached the villages, the \textbf{land holding of the Lingayats (the dominant caste) increased substantially, while that of the Adi Karnataka (scheduled caste) decreased}.
\item In the \textbf{wet (irrigated) areas}, outsiders also began \textbf{buying land} (investing in the land market), and \textbf{most of the benefit of irrigation was appropriated by the big (Lingayat) farmers} --- who controlled the cooperatives (loans at 6\%, half the bank rate) and could \textbf{buy fertilisers/chemicals on credit} (paying after harvest), options the small and marginal farmers lacked. So the \textbf{inequality between the dominant caste (Lingayats) and the scheduled castes (Adi Karnataka) increased}.
\end{itemize}
\end{KeyConcept}

\begin{QuickRevision}
\begin{itemize}
\item Scarlett Epstein
\item longitudinal study
\item Mandya district, Karnataka
\item wet (irrigated) areas
\item buying land
\item buy fertilisers/chemicals on credit
\end{itemize}
\end{QuickRevision}

\section{Mechanisation and the Labourers: Dhanagare}

\begin{KeyConcept}
\begin{itemize}
\item \textbf{D. N. Dhanagare}: mechanisation \textbf{improved the farmers' condition but was catastrophic for labourers} --- mostly \textbf{migrant labourers from Bihar and UP} (cyclical migration, staying for the season).
\item The costs: in \textbf{12,000 villages of Punjab there were 2 lakh farm machines (1984)}; \textbf{5,000 farm-accident deaths in 1978}; about \textbf{1,000 farm workers lost limbs in one year (1985)} in Punjab/Haryana/Western UP; and by 1985 \textbf{over 10,000 migrant labourers had died}. Agricultural labourers were \textbf{not covered by the Workmen's Compensation Act} (they were not counted as 'workers'), so they got no compensation.
\item \textbf{Pesticide poisoning}: by 1984 there were \textbf{7.5 lakh acute poisoning cases} (International Development Research Centre, Ottawa), of which \textbf{one-third were from India} --- and the victims were agricultural labourers, rarely the rich farmers.
\item As labourers fled to industrial towns (for wages and compensation), the Green Revolution areas faced a labour shortage and created a \textbf{new form of bondage}: landowners gave labourers land and brought their families --- the \textbf{whole family became servants of the landowner}.
\end{itemize}
\end{KeyConcept}

\begin{QuickRevision}
\begin{itemize}
\item Dhanagare: mechanisation helped farmers, catastrophic for migrant labourers (Bihar/UP, cyclical migration).
\item 2 lakh machines in 12,000 Punjab villages (1984); 5,000 deaths (1978); 1,000 lost limbs (1985); 10,000 migrant deaths by 1985; no compensation (not under Workmen's Compensation Act).
\item 7.5 lakh pesticide poisonings (1/3 from India; victims = labourers); new bondage: whole family = servants of the landowner.
\end{itemize}
\end{QuickRevision}

\section{Environment, Decline and New Politics}

\begin{KeyConcept}
\begin{itemize}
\item \textbf{Environment}: Punjab and parts of Haryana are \textbf{cancer-prone} --- arsenic and uranium in the groundwater; Bihar's groundwater too has arsenic (journal \emph{Water}, from illegal sand-mining), but the poor cannot afford RO systems --- so cancer centres have also appeared in Bihar.
\item \textbf{Decline}: the productivity of Green Revolution areas has \textbf{declined since the 1990s} --- and this has produced \textbf{new forms of politics} --- the \textbf{Jat reservation} demand (Punjab/Haryana), the \textbf{Patidar} agitation (Gujarat), and similar movements in Andhra.
\item Why reservation? The Green Revolution societies \textbf{never invested in modern institutions} (schools etc.) and remained conservative; they cannot compete in private jobs (which need communication skills), but \textbf{government jobs are a level playing field} --- hence the reservation agitation.
\end{itemize}
\end{KeyConcept}

\begin{QuickRevision}
\begin{itemize}
\item Environment: arsenic/uranium groundwater (Punjab/Haryana cancer-prone; Bihar arsenic).
\item Productivity declined since the 1990s $\rightarrow$ new politics (Jat reservation, Patidar agitation); no investment in modern institutions $\rightarrow$ private jobs inaccessible, government jobs = level playing field.
\end{itemize}
\end{QuickRevision}
